\documentclass[11pt,a4paper]{article}

\usepackage[margin=2.2cm]{geometry}
\usepackage{booktabs}
\usepackage{adjustbox}
\usepackage{enumitem}
\usepackage{microtype}
\usepackage[hidelinks]{hyperref}

\setlength{\parindent}{0pt}
\setlength{\parskip}{0.55em}
\setlist[itemize]{leftmargin=1.5em,itemsep=0.2em,topsep=0.2em}

\title{Class-Weight Sensitivity Analysis for 30-Day Hospital Readmission Prediction}
\author{Xiaohan Mu}
\date{}

\begin{document}
\maketitle

\section{Purpose}

This report summarises a focused class-weight sensitivity analysis for the four main tuned models: Regularised Logistic Regression, Decision Tree, Random Forest, and XGBoost. The purpose of the experiment was to investigate whether giving greater importance to the minority 30-day readmission class could improve discrimination or reduce the false-positive burden required to maintain high recall.

The readmission outcome is strongly imbalanced, with approximately 9\% of encounters belonging to the positive class. Class weighting is therefore a plausible method for making errors on readmitted patients more costly during model fitting. However, weighting does not add new predictive information to the dataset. Its value depends on whether changing the training objective leads the model to rank readmitted patients more effectively or only shifts the scale of its predicted probabilities.

\section{Experimental design}

The sensitivity analysis reused the same 18 predictors and the same 60/20/20 model-training, validation, and test split used in the main optimisation workflow. For each model, the previously selected non-weight hyperparameters were held fixed while the positive-class weight was varied.

A common manual sequence of positive-to-negative class-weight ratios was tested:

\[
1,\;2,\;4,\;6,\;8,\;10.
\]

Automatic reference settings such as \texttt{balanced} and, for Random Forest, \texttt{balanced\_subsample}, were also retained where supported.

For each weight setting, five-fold stratified cross-validation was used on the model-training data to calculate mean AUPRC, AUROC, and Brier score. AUPRC was treated as the primary discrimination metric because readmission is the minority class. The fitted configuration was then applied to the validation set, where a threshold was selected to achieve at least 80\% recall while minimising the false-positive rate (FPR).

The experiment was interpreted as a sensitivity analysis rather than as a complete replacement for the earlier joint hyperparameter searches. This is important because class weight can interact with other hyperparameters such as regularisation strength, tree structure, and boosting settings.

\section{Results}

\subsection{Regularised Logistic Regression}

The expanded Logistic Regression tuning now includes the intermediate class-weight values and selected a positive-class weight of 2, represented as \texttt{\{0:1, 1:2\}}. The final selected additive model retained $C=0.01$ and L2 regularisation. Its cross-validation AUPRC was approximately 0.1516 and its validation-selected threshold increased from approximately 0.067 in the earlier unweighted model to 0.126.

The class-weight sensitivity experiment showed that discrimination was almost unchanged across the tested weights. The mean cross-validation AUPRC was approximately 0.15148 at weight 1 and 0.15155 at weight 2, a difference of only around 0.00008. Higher weights did not provide a clear further improvement. The corresponding AUROC values were also very similar.

The effect on the probability scale was much larger than the effect on discrimination. Increasing the weight shifted predicted probabilities upward and therefore required a progressively higher classification threshold to retain approximately the same recall. Despite this large change in threshold, the false-positive rate remained around the same general level once the threshold was re-selected for at least 80\% recall.

The final weight-2 Logistic Regression achieved test AUPRC 0.1497, AUROC 0.6368, recall 0.8003, FPR 0.6362, and Brier score 0.0860. The Brier score is worse than the approximately 0.0800 obtained by the earlier unweighted version, while discrimination is almost unchanged. Therefore, the selected weight of 2 should be interpreted as a mild shift in the training objective rather than evidence that stronger class weighting substantially improves the prediction problem.

\subsection{Decision Tree}

When the already selected tree structure was held fixed, the unweighted version produced the numerically highest cross-validation AUPRC, approximately 0.1401, compared with approximately 0.1398 for the balanced reference. The differences in discrimination across weights were small.

The stronger pattern occurred in probability quality. The validation Brier score was approximately 0.081 for the unweighted tree but increased to approximately 0.237 under balanced weighting. Stronger positive weights progressively increased the Brier score without producing a corresponding increase in AUPRC.

This does not necessarily contradict the earlier broad hyperparameter search, which selected a balanced Decision Tree. The earlier search changed class weight together with tree depth, leaf size, feature sampling, and other structural settings. The sensitivity experiment instead asks what happens when class weight alone changes while the selected tree structure remains fixed. The two analyses therefore answer different questions.

The result suggests that balanced weighting is not independently responsible for a clear increase in Decision Tree discrimination and can substantially distort the raw probability estimates. The Decision Tree also continued to have a high false-positive burden at the required recall level, so this finding does not make it a stronger final-model candidate.

\subsection{Random Forest}

Random Forest gave the clearest result. The unweighted model achieved the highest mean cross-validation AUPRC, approximately 0.1521. AUPRC gradually decreased as the positive-class weight increased, while AUROC remained similar or declined slightly.

Probability quality deteriorated much more strongly. The validation Brier score increased from approximately 0.080 at weight 1 to approximately 0.23 at weights close to the full class-imbalance ratio. Small changes in class weight produced only very small differences in FPR at the validation threshold, while larger weights provided no meaningful discrimination benefit.

This strongly supports the original Random Forest selection of \texttt{class\_weight=None}. The result indicates that the forest already extracts the available ranking information from the minority class without requiring additional weighting. Increasing the training penalty on positive-class errors does not improve separation and mainly shifts the probability estimates away from the observed outcome frequency.

\subsection{XGBoost}

XGBoost showed a less monotonic pattern. In the fixed-structure sensitivity experiment, the numerically highest mean cross-validation AUPRC occurred at a positive-class weight of 10, at approximately 0.1519, while weights between 2 and 8 produced values very close to this result. The original joint hyperparameter search selected \texttt{scale\_pos\_weight=6}.

The numerical differences in AUPRC were small relative to the variability between cross-validation folds, so the result does not provide strong evidence that weight 10 is meaningfully superior to weight 6 or to the moderate-weight configurations. The operating-point results were also close.

As with the other models, the clearest effect of strong weighting was deterioration in probability quality. Validation Brier score increased from approximately 0.080 at weight 1 to approximately 0.23 at weight 10. The existing weight-6 tuned model also has a substantially higher Brier score than the unweighted Random Forest. The XGBoost result therefore suggests that some positive weighting may be useful to the boosting optimisation objective, but stronger weighting does not clearly solve the underlying discrimination problem and makes the raw probability estimates less suitable for direct interpretation as absolute readmission risks.

\section{What the different selected weights mean}

The experiment does not show that one class-weight value should be optimal for every model. Class weighting interacts with the way each algorithm learns.

Regularised Logistic Regression selected a mild weight of 2, Random Forest preferred no additional weighting, the main Decision Tree tuning selected balanced weighting, and XGBoost selected a stronger weight of 6 during the joint hyperparameter search. This variation indicates that class weight is model-dependent.

A weight changes the loss or splitting objective used during training, and different algorithms respond to that change in different ways. For a linear model, weighting changes the relative cost of positive and negative errors and therefore shifts the fitted decision boundary. In tree-based methods it can change which splits or leaves are considered useful. In boosting, it changes how strongly positive-class errors influence successive boosting updates.

Therefore, a stronger selected weight should not be interpreted as meaning that a model has discovered more useful information in the readmission class, nor should a weight of 1 be interpreted as meaning that the minority class is being ignored. The useful question is whether the chosen weight improves the evaluation objective.

In this project, ``best'' primarily means the configuration that gives the strongest cross-validated AUPRC during model selection, with AUROC and probability-quality metrics used as supporting evidence. After the model is selected, the validation threshold is chosen separately to meet the 80\% recall requirement while minimising FPR. A class weight may therefore be numerically selected because it produces a small AUPRC improvement even if another weight has a better Brier score or slightly better FPR.

The sensitivity results also show why a purely numerical winner should not automatically be treated as a meaningful winner. Several AUPRC differences were extremely small compared with fold-to-fold cross-validation variation. These settings should be regarded as practically similar unless the improvement is consistent across other metrics or clearly larger than the observed variability.

\section{Overall interpretation}

The main finding is that class weighting did not substantially improve the models' ability to separate readmitted from non-readmitted patients. Across the tested weight ranges, AUPRC and AUROC changed only slightly. At the validation thresholds chosen to maintain at least 80\% recall, the false-positive rates also remained high and changed by only small amounts.

The much stronger effect was on the predicted probability scale. Larger positive-class weights caused the models to assign higher probabilities to patients in general. The threshold therefore also had to increase to maintain a comparable operating point. Because the ordering of patients changed much less than the probability values themselves, AUPRC and AUROC remained almost stable while Brier scores often deteriorated substantially.

This means that class weighting and threshold adjustment partly address the same practical problem from different directions. Class weighting makes missing positive cases more expensive during training, while threshold adjustment makes it easier for an already trained model to classify a patient as positive. If weighting does not meaningfully improve the ranking of patients, an unweighted model with a lower threshold can behave similarly to a weighted model with a higher threshold.

The results therefore suggest that the high false-positive burden is not mainly caused by the models failing to pay enough attention to the minority readmission class. Instead, the available predictors appear to provide only limited separation between patients who will and will not be readmitted. Increasing class weight changes how aggressively models respond to positive cases, but does not create additional predictive information.

\section{Implications for final model selection}

The class-weight experiment should be treated as a robustness and interpretation analysis rather than as a reason to replace every model with the numerical weight winner from the fixed-structure sensitivity test.

The expanded Logistic Regression search now legitimately selects weight 2 as part of its joint tuning process. Random Forest has strong evidence supporting the unweighted configuration. For Decision Tree and XGBoost, the differences between weight settings should be interpreted in the context of their full joint hyperparameter searches and their calibration behaviour rather than using the sensitivity AUPRC alone.

For the final model comparison, class weight should therefore be considered together with threshold-independent discrimination, performance at the selected high-recall operating point, calibration and probability quality, and model complexity and interpretability.

\section{Conclusion}

The class-weight sensitivity analysis showed that different algorithms can prefer different levels of positive-class weighting, but no tested weighting strategy produced a substantial improvement in readmission discrimination. Logistic Regression selected a mild positive weight of 2 after its expanded tuning grid, while Random Forest continued to favour no additional weighting. Decision Tree and XGBoost could select stronger weights in their joint searches, but the sensitivity analysis showed that the discrimination differences between many weight settings were small.

The most consistent effect of stronger weighting was an upward shift in predicted probabilities and a deterioration in Brier score. After thresholds were independently adjusted to maintain approximately 80\% recall, the false-positive burden remained broadly similar. The results therefore suggest that class weighting can modify the model's emphasis on minority-class errors but cannot overcome the limited separation available from the current predictors.

Class weighting should therefore be viewed as one component of model optimisation rather than a solution to the readmission imbalance problem. The next evaluation stage should focus on calibration, interpretation of the strongest candidate models, and subgroup performance rather than further expansion of the class-weight search.

\end{document}
